\chapter{How a finite calculation controls every conformal weight}
\label{ch:finite_closure}

The definition of $U$ uses infinitely many possible iterated mode
products. A calculation through weight four cannot, by itself, say
what happens at weight twenty. The purpose of this chapter is to bridge
that gap. We first prove an abstract recursive criterion, then explain
how the particular finite spaces in the Haagerup construction satisfy
its hypotheses. The central idea is that higher generators come with
actual recipes in lower generators. Those recipes let Jacobi identities
replace an unbounded search by a terminating induction.

The argument is useful far beyond this example. It also teaches how
to read a computer-assisted proof: exact ranks, exact memberships,
ambient dimensions, and recursive provenance support different logical
steps. None can be substituted for another merely because the final
numbers agree.

\section{Generation and strong generation}
\index{strong generation}

For this chapter write $D=L_{-1}$ for translation and
$:ab:=a_{(-1)}b$. We retain the mode convention
$Y(a,z)=\sum_{n\in\Z}a_{(n)}z^{-n-1}$. The weight rule is
\begin{equation}\label{eq:c-mode-weight}
 \wt(a_{(n)}b)=\wt(a)+\wt(b)-n-1.
\end{equation}
Products with $n\geq0$ are the singular OPE coefficients. Products
with negative $n$ are normal products and their derivatives.

\begin{definition}
For a finite-dimensional homogeneous positive-weight space $S$, its
\emph{differential normal span} $\mathcal N(S)$ is the span of the
vacuum and all expressions built from $S$ using $D$ and arbitrarily
parenthesized normal products. The space $S$ strongly generates a
vertex algebra when this span is the whole algebra.
\end{definition}

Allowing arbitrary parentheses avoids an unnecessary preliminary
rewriting theorem. Normal multiplication in a VOA is neither literally
commutative nor literally associative. Its defects are controlled by
singular products and derivatives. If a preferred ordering is desired
for a computer representation, its validity must be justified by those
identities.

Vertex generation permits all modes at every step. Strong generation
uses only normal products and derivatives after the strong generators
have been named. The distinction is visible immediately in weight two.
If $a,b$ both have weight two, then $a_{(0)}b$ has weight three.
It need not be a normal expression in $a,b$: the smallest normal
product of them has weight four. Such a weight-three field may need to
be added as a strong generator even though it is already a vertex
product of the original generators.

\begin{example}[A one-current test]
A Heisenberg current $J$ satisfies $J_{(1)}J=\one$ and
$J_{(0)}J=0$, with all higher products zero. These singular products
already lie in $\mathcal N(J)$. Wick reduction then shows that every
singular product of two normal expressions in $J$ lies there as well.
Thus one current strongly generates its Heisenberg VOA. The point is
the verified OPE closure, not merely the observation that the first
few Fock-space dimensions agree with partition numbers.
\end{example}

\section{The two identities that drive the proof}

Translation obeys sesquilinearity, and zero modes are derivations:
\[
 [D,a_{(n)}]=-n a_{(n-1)},\qquad
 a_{(0)}:xy:=:(a_{(0)}x)y:+:x(a_{(0)}y):.
\]
The second identity is especially valuable because it has no correction
term. Also $[D,a_{(0)}]=0$.

For arbitrary modes the efficient language is the $\lambda$-bracket,
\[
 [a_\lambda b]=\sum_{n\geq0}\frac{\lambda^n}{n!}a_{(n)}b.
\]
The noncommutative Wick identity is
\begin{equation}\label{eq:c-wick}
 [a_\lambda:bc:]=:[a_\lambda b]c:+:b[a_\lambda c]:
          +\int_0^\lambda[[a_\lambda b]_\mu c]\,d\mu.
\end{equation}
The integral is a formal polynomial integral, not an analytic operation.
It accounts for the double contractions that a naive Leibniz rule
would miss.

\begin{lemma}[Filtered Wick reduction]\label{lem:c-filtered-wick}
Suppose singular products of generator pairs of total weight less than
$w$ belong to $\mathcal N(S)$. The same holds for normal expressions
of total input weight less than $w$. At weight $w$, it holds whenever
one input is a derivative or a normal product of positive-weight
expressions.
\end{lemma}
\begin{proof}
Induct on total input weight. Sesquilinearity removes a derivative and
reduces that weight. Skew symmetry moves a composite to the right.
Apply \eqref{eq:c-wick}. In each of the first two terms one factor
has been omitted from the bracket, so its total input weight is smaller.
For the integral term, the coefficient $a_{(j)}b$ has weight
$\wt(a)+\wt(b)-j-1$. Pairing it with $c$ gives total weight at
most $\wt(a)+\wt(b)+\wt(c)-1$, again smaller. If the coefficient
is scalar its bracket vanishes. Normal multiplication and derivatives
preserve the span by definition. This proves the induction without
assuming that normal words are linearly independent.
\end{proof}

Notice what this lemma does not do. It reduces products containing
composites, but it does not settle a new pair of primitive strong
generators. We need one more source of structure for those pairs.

\section{A recursive strong-generation theorem}

Suppose we have finite spaces $S_d\subset V_d$ for $1\leq d\leq h$,
with $h\geq2$. Let $K=S_1$, $A=S_2$, and
$\mathcal N=\mathcal N(S_1+\cdots+S_h)$. A subscript $d$ on
$\mathcal N$ denotes its homogeneous part of conformal weight $d$.
Here $K$ denotes the
finite current space; in the Haagerup application it is the weight-one
part of the affine algebra called $K$ in the preceding chapter.
A higher field has
\emph{recursive provenance} if, modulo lower normal expressions, it is
a zero product of a weight-two field with a field one weight lower.

\begin{theorem}[Finite recursive closure]\label{thm:c-recursive}
\index{finite closure theorem}
Assume $V$ is nonnegatively integer graded with $V_0=\C\one$, and
that the following identities hold in $V$:
\begin{enumerate}
\item For $k\in K$, $a\in A$, $s\in S_d$, and every $n\geq0$,
both $k_{(n)}s$ and $a_{(n)}s$ belong to $\mathcal N$.
\item For $3\leq d\leq h$,
\[
 S_d\subset A_{(0)}S_{d-1}
       +\mathcal N_d(S_1+\cdots+S_{d-1}).
\]
\end{enumerate}
Then $\mathcal N$ is a vertex subalgebra. Hence the $S_d$ strongly
generate the vertex algebra that they generate.
\end{theorem}
\begin{proof}
First establish $a_{(0)}\mathcal N\subset\mathcal N$ for $a\in A$.
This follows from the first hypothesis on generators, the zero-mode
derivation identity, and commutation with $D$. It uses no unproved
closure of other singular products.

Now order generator pairs $(s,t)$ lexicographically by
\[
 (\wt(s)+\wt(t),\wt(s)).
\]
We prove singular closure in that order. If the left weight is one or
two, it is a hypothesis. Otherwise replace $s$ by terms $a_{(0)}b$
and lower normal expressions, where $\wt(b)=\wt(s)-1$. The normal
terms are controlled by Lemma~\ref{lem:c-filtered-wick}. Jacobi gives
\begin{equation}\label{eq:c-recursive-jacobi}
 (a_{(0)}b)_{(n)}t
     =a_{(0)}(b_{(n)}t)-b_{(n)}(a_{(0)}t).
\end{equation}
In the first term the inner pair has lower total weight, and the
outer $a_{(0)}$ preserves $\mathcal N$.

The first hypothesis gives a normal expression for $a_{(0)}t$.
Every composite or derivative in that expression reduces the second
term by filtered Wick. The only other possibility is a single
undifferentiated generator of weight $\wt(t)+1$. Pairing that
generator with $b$ has the original total weight but smaller left
weight, so it is earlier in the lexicographic order. If $\wt(t)=h$
there is no such generator. The induction therefore terminates.

All generator singular products now close. Filtered Wick extends this
to all normal expressions. For negative modes use
\[
 x_{(-r)}y=\frac1{(r-1)!}:(D^{r-1}x)y:,
 \qquad r\geq1.
\]
Thus every mode preserves $\mathcal N$, proving the assertion.
\end{proof}

The strength of the theorem is its finite range. The largest output
weight in the first hypothesis is $h+1$, arising from a zero product
of weights two and $h$. Directly checking every pair of strong
generators would reach weight $2h-1$. More importantly, the theorem
explains why those finitely many checks suffice in all weights.

\begin{remark}[Why provenance is indispensable]
Suppose an arbitrary extra vector were added at weight four because
it made a character coefficient look correct. Its brackets with
another weight-four vector would not be controlled by the theorem.
The recipe $S_4\subset A_{(0)}S_3+$ lower normals is what makes
\eqref{eq:c-recursive-jacobi} available. A dimension is not a recipe.
\end{remark}

\section{The finite problem in the Haagerup algebra}

Return to the actual $U\subset B$. Let $K$ be its six known currents,
and let $S_2,S_3$ be the constructed spaces of dimensions $120,956$.
Their lower normal space is
\[
 N_3=K_{(-1)}S_2+DS_2,\qquad \dim N_3=634,
\]
and the exact provenance statement is
\[
 S_3=N_3+\Span\{a_{(0)}b:a,b\in S_2\}.
\]
Choose actual lifts of the quotient $S_3/N_3$ and an actual current
triplet $P_4$ of highest weight $(0,2)$. Here highest weights, rather
than representation dimensions, are being used: $(0,2)$ means the
representation of dimensions $(1,3)$.

Set
\begin{align}
 N_4&=K_{(-1)}S_3+DS_3+(S_2)_{(-1)}S_2,
 &W_4&=N_4+P_4,\\
 N_5&=K_{(-1)}W_4+DW_4+(S_2)_{(-1)}S_3.
 \label{eq:c-normal-spaces}
\end{align}
These are the full displayed spans; their exhaustion of the relevant
normal expressions is justified below. A solver may use selected
columns that lie in them, but those selected columns must not silently
replace the full spaces in an algebraic argument.

\begin{constructioninput}[Lower closure and provenance]
The actual spaces are stable under both current zero-mode algebras.
All singular products of currents and weight-two fields with $S_2$
belong to the expected spaces in weights at most three. The triplet
is constructed from an actual $a_{(0)}p$, with $a\in S_2$ and
$p\in S_3$, by current projectors and lowering. It satisfies
\[
 P_4\subset(S_2)_{(0)}S_3+N_4,
 \qquad L_1P_4\subset S_3,\quad L_2P_4\subset S_2.
\]
The stated inclusions are exact vector identities in the ambient
lattice algebra and are available before the next two gates.
\end{constructioninput}

The two remaining central memberships are
\begin{align}
 (S_2)_{(0)}S_3&\subset W_4,
 \label{eq:c-gate4}\\
 (S_2)_{(0)}P_4&\subset N_5.
 \label{eq:c-gate5}
\end{align}
We call them the weight-four and weight-five gates. The name refers
to the output weight, not an empirical test that can be replaced by
sampling a few products.

\section{Raising operators turn seeds into a complete proof}

In this representation-theoretic section, $E_L,E_R,F_L,F_R$
abbreviate the zero-mode operators of the corresponding current states.
Let $P=S_3/N_3$, of dimension $322$. For a fixed highest field
$a\in S_2$, define
\[
 f_a:P\longrightarrow B_4/W_4,\qquad[p]\longmapsto[a_{(0)}p].
\]
Lower closure makes the map well defined. Because $a$ is highest,
its kernel is stable under the two commuting raising operators
$E_L,E_R$.

\begin{lemma}[Finite raising saturation]\label{lem:c-raising}
\index{raising operators!finite saturation}
Suppose $C\subset\ker f_a$ and
$P=E_LP+E_RP+C$. Then $f_a=0$.
\end{lemma}
\begin{proof}
On $P/\ker f_a$ the equation becomes $M=E_LM+E_RM$.
The ideal generated by the two raising operators is nilpotent because
the finite-dimensional representation has bounded weights. Iterating
gives $M=(E_L,E_R)^rM=0$ for sufficiently large $r$.
\end{proof}

\begin{example}[One lowest vector per irreducible]
In an irreducible $\mathfrak{sl}_2$ representation of dimension $m$,
the raising image has codimension one. Its missing direction is the
lowest-weight vector. In a product of two irreducibles,
$E_LP+E_RP$ again has codimension one: only the joint-lowest vector
is missing. Therefore a direct sum of $r$ irreducible constituents,
counted with multiplicity, needs $r$ such seeds. This explains the
size of the Haagerup gate before any large matrix is written.
\end{example}

The exact current decomposition of $P$, in representation dimensions,
is
\begin{align*}
 P={}&3(1,3)+(1,5)+5(1,7)+2(1,9)+3(1,11)+(1,13)\\
 &+2(1,15)+(1,19)+(2,2)+(2,4)+2(2,6)+2(2,8)\\
 &+2(2,10)+(2,12)+(2,14).
\end{align*}
There are $28$ constituents and total dimension $322$. Hence the
raising image has dimension $294$. The nine primitive tests are the
eight listed nontrivial weight-two multiplets and one independent neutral
stress; affine normal states and the ambient stress are already controlled.
For each of these highest fields in the first input, exact memberships
on a spanning set of the
$28$ missing directions prove the complete map vanishes. This is the
conceptual content of the weight-four computation.

Some residual outputs have weight $(0,0)$ and require a little care.
A vector being a highest class modulo $W_4$ does not mean its chosen
ambient representative is literally a singlet. The following projection
argument recovers the full membership.

\begin{lemma}[Recovering a residual singlet]
Let $v\in B_4$ have current weight $(0,0)$ and satisfy
$E_Lv,E_Rv\in W_4$. Let $\Pi$ be the polynomial projector onto
joint singlets. If $\Pi v$ lies in $N_4$, then $v\in W_4$.
\end{lemma}
\begin{proof}
Each nonconstant factor in the projector contains $FE$, so acts as
zero on the class of $v$ modulo $W_4$. Thus $[\Pi v]=[v]$.
The hypothesis makes the left side zero.
\end{proof}

There is no mysterious infinite representation theory in this
projector. At ambient weight $d$, a momentum has norm at most
$\sqrt{2d}$, so a Cartan eigenvalue obeys
$m^2\leq2d(h,h)$. At weight four the possible highest weights are
at most $4$ and $24$ for the two current factors. On weight zero,
$FE$ acts in highest weight $j$ by $j(j+2)/4$. Therefore
\[
 \Pi=\prod_{j=2,4}\left(1-\frac{4F_LE_L}{j(j+2)}\right)
 \prod_{j=2,4,\ldots,24}
                 \left(1-\frac{4F_RE_R}{j(j+2)}\right)
\]
is an exact finite polynomial. Exact projected residuals complete the
exceptional seed memberships.

\section{Why the weight-five gate is a Jacobi problem}

After the weight-four gate, Virasoro commutators recover the positive
products that will be needed next. This step should be proved before
any use of the weight-five gate, since its normal-space argument
depends on precisely these products.

\begin{lemma}[The intermediate positive-mode controls]
\label{lem:c-positive-bootstrap}
The lower products, the direct $L_1P_4,L_2P_4$ controls, and the
weight-four gate imply
\begin{align*}
 (S_2)_{(1)}S_3&\subset S_3,&
 (S_2)_{(2)}S_3&\subset S_2,&
 (S_2)_{(3)}S_3&\subset K,\\
 K_{(1)}W_4&\subset S_3,&
 K_{(2)}W_4&\subset S_2,&
 K_{(3)}W_4&\subset K,\\
 L_1N_5&\subset W_4,&
 L_2N_5&\subset S_3.&&
\end{align*}
The next mode in each of the first two lines is scalar; further
modes vanish by weight.
\end{lemma}
\begin{proof}
Current products with $S_3$ follow first from its provenance.
For example, if $k\in K$ and $a,b\in S_2$, then
\[
 k_{(1)}(a_{(0)}b)
  =(k_{(0)}a)_{(1)}b+(k_{(1)}a)_{(0)}b
                           +a_{(0)}(k_{(1)}b).
\]
Every right-hand term belongs to $S_2$ by lower closure.
Higher current modes and normal inputs are handled by the same
commutator and iterate identities.

For a general weight-two vector, the needed Virasoro commutators are
\begin{align*}
 [L_1,a_{(n)}]&=(2-n)a_{(n+1)}+(L_1a)_{(n)},\\
 [L_2,a_{(n)}]&=(3-n)a_{(n+2)}
                     +3(L_1a)_{(n+1)}+(L_2a)_{(n)}.
\end{align*}
These give $L_1S_3\subset S_2$, $L_2S_3\subset K$ and,
on the three defining summands of $N_4$,
$L_1N_4\subset S_3$, $L_2N_4\subset S_2$.
For $L_2(a_{(-1)}b)$ the term $(L_2a)_{(-1)}b$
must be retained: $L_2a$ is scalar, but its $(-1)$-mode
acts by that scalar. The direct controls on $P_4$ extend these
inclusions to $W_4$ before the weight-five gate.

The primitive weight-two multiplets and the two neutral singlets
are quasiprimary. Indeed $L_1$ is current-equivariant, and none
of their current types occurs in $K$, whose types are $(3,1)$
and $(1,3)$. Thus $L_1a=0$ for these $a$. For $p\in S_3$,
the weight-four gate now gives
\begin{align*}
 2a_{(1)}p&=L_1(a_{(0)}p)-a_{(0)}L_1p\in S_3,\\
 a_{(2)}p&=L_1(a_{(1)}p)-a_{(1)}L_1p\in S_2,\\
 2a_{(3)}p&=L_2(a_{(1)}p)-a_{(1)}L_2p\in K.
\end{align*}
In the last line $(L_2a)_{(1)}p=0$, because the positive
modes of the vacuum are zero. This differs from the normal
product term just discussed. The current-normal and derivative
parts of $S_2$ use the general commutators and current iterates,
so the first line of the lemma holds on all of $S_2$.

Next use $P_4\subset(S_2)_{(0)}S_3+N_4$ in the displayed
current identity, now allowing $b\in S_3$. It gives
$K_{(1)}P_4\subset S_3$. Since a current is primary,
$[L_1,k_{(1)}]=-k_{(2)}$ and
$[L_2,k_{(1)}]=-k_{(3)}$.
The established Virasoro controls therefore give the remaining
current inclusions on $W_4$.

Finally apply $L_1,L_2$ to the three summands of $N_5$.
The current and derivative summands use the controls already proved.
For the normal input of weights two and three,
\begin{align*}
 L_1(a_{(-1)}p)
   &=3a_{(0)}p+(L_1a)_{(-1)}p+a_{(-1)}L_1p,\\
 L_2(a_{(-1)}p)
   &=4a_{(1)}p+3(L_1a)_{(0)}p
                  +(L_2a)_{(-1)}p+a_{(-1)}L_2p.
\end{align*}
The respective outputs lie in $W_4$ and $S_3$. Reordering a
normal term uses only lower products that have already been
established. This completes the proof.
\end{proof}

The point of spelling out this intermediate lemma is to make the
logical order visible. We now possess the positive products needed
to analyze all of the normal subspace $N_5$; no calculation at
weight five has been used to establish its own hypotheses.

At this stage every normal word of weight five reduces to $N_5$.
A leading derivative reduces to $W_4$; a normal product splits into
weights $1+4$ or $2+3$, after reversing the factors if necessary.
The reordering corrections are derivatives of $K_{(n)}W_4$ or
$(S_2)_{(n)}S_3$, which have just been controlled. Higher negative
modes reduce by translation:
\[
 x_{(-r)}v=\frac{D(x_{(-(r-1))}v)-x_{(-(r-1))}Dv}{r-1},
 \qquad r\geq2.
\]
Thus no unproved $S_3$-by-$S_3$ singular closure is hidden in the
following normal-space containments.

One obtains, before the final gate,
\[
 (S_2)_{(0)}N_4\subset N_5,
 \qquad (N_3)_{(0)}S_3\subset N_5.
\]
The second inclusion follows from the iterate identity
\[
 (:xa:)_{(0)}p=
 \sum_{i\geq0}x_{(-1-i)}(a_{(i)}p)
 +\sum_{i\geq0}a_{(-1-i)}(x_{(i)}p),
\]
for a current $x$. Every inner mode is already controlled, and grading
makes the sums finite.

If $P_4\subset N_4$, nothing remains. Otherwise $Q=W_4/N_4$ is the
current triplet. Define
\[
 \beta(a,p)=[a_{(0)}p]\in Q,
 \qquad F(a,q)=[a_{(0)}\widetilde q]\in B_5/N_5.
\]
The normal-space controls make $F$ independent of the lift.

\begin{lemma}[Jacobi exchanges the difficult input]
If $a_{(0)}b\in N_3$, then
\begin{equation}\label{eq:c-bottom-jacobi}
 F(a,\beta(b,p))=F(b,\beta(a,p)).
\end{equation}
\end{lemma}
\begin{proof}
The difference is the class of $(a_{(0)}b)_{(0)}p$, which lies
in $N_5$ by the preceding inclusion.
\end{proof}

Here \emph{charged} means a nontrivial representation of the retained
current zero modes, while \emph{neutral} means a current singlet.
Write $q_-$ for a lowest-weight vector of the triplet $Q$.
The finite construction supplies two nonzero bottom witnesses
$\beta(b,p)=cq_-$, $c\ne0$. For each charged highest field $a$,
one chooses such a witness for which $a_{(0)}b\in N_3$ and the
weight of $a_{(0)}p$ is outside the three weights of $Q$. Then
$\beta(a,p)=0$, so \eqref{eq:c-bottom-jacobi} gives
$F(a,q_-)=0$. Raising this \emph{proved zero identity} gives the
other two triplet inputs; current covariance supplies the descendants
of $a$.

The emphasized order matters. The weight exclusion holds at the chosen
bottom input. It need not hold after that input is raised. What is
raised is the zero identity already established by Jacobi, not an
unverified assertion that all later intermediate weights are absent.

The remaining neutral family is reduced to charged brackets. Their
current span exhausts $P$ by the same nilpotent-ideal argument, now
using lowering. Here is one explicit reduction. Let $e$ be the
retained neutral singlet and let $a_{\rm ch},p_{\rm bot}$ be a
charged bottom witness with
$\beta(a_{\rm ch},p_{\rm bot})=c q_-$, $c\ne0$.
Charged spanning gives
\[
 e_{(0)}a_{\rm ch}
       =\sum_i t_i(x_i)_{(0)}y_i+n,\qquad n\in N_3,
\]
with charged $x_i,y_i\in S_2$. Applying zero modes to
$p_{\rm bot}$ gives
\[
 (e_{(0)}a_{\rm ch})_{(0)}p_{\rm bot}
 =\sum_i t_i[(x_i)_{(0)},(y_i)_{(0)}]p_{\rm bot}
                            +n_{(0)}p_{\rm bot}\in N_5.
\]
The inner charged products belong to $W_4$ by the weight-four
gate. The next charged products are controlled on both $N_4$
and $P_4$ by the normal-space results and the charged portion
of the weight-five argument. A second Jacobi identity yields
\[
 e_{(0)}((a_{\rm ch})_{(0)}p_{\rm bot})
  =(a_{\rm ch})_{(0)}(e_{(0)}p_{\rm bot})
                    +(e_{(0)}a_{\rm ch})_{(0)}p_{\rm bot}\in N_5.
\]
The nonzero coefficient $c$ proves the neutral bottom identity.
The neutral zero mode commutes with current zero modes, so covariance
gives the whole triplet. The other neutral stress follows by
subtracting the retained stress and the known Sugawara stresses
from the ambient stress.

\begin{constructioninput}[The two completed gates]
For the actual finite spaces, the raising-seed memberships and their
singlet residuals establish \eqref{eq:c-gate4}. The nonzero bottom
witnesses, lower bracket memberships, and charged spanning identities
described above establish \eqref{eq:c-gate5}. All equalities are over
$\Q(\sqrt{13},\sqrt{-3})$ in the fixed lattice basis. The gates are
finite construction inputs; the recursive theorem is their
all-conformal-weight consequence.
\end{constructioninput}
The construction manuscript's certificate ledger lists the completed
lower-closure, actual-triplet, weight-four, charged-span, and
weight-five records. The source guide gives the manuscript location;
the finite identities remain inputs here, rather than calculations
independently replayed by this textbook.

\section{Properness is now an all-weight conclusion}

We can now close the last positive modes explicitly. For a
quasiprimary primitive $a\in S_2$ and $p\in P_4$, the
weight-five gate and Lemma~\ref{lem:c-positive-bootstrap} give
\begin{align*}
 2a_{(1)}p&=L_1(a_{(0)}p)-a_{(0)}L_1p\in W_4,\\
 a_{(2)}p&=L_1(a_{(1)}p)-a_{(1)}L_1p\in S_3,\\
 2a_{(3)}p&=L_2(a_{(1)}p)-a_{(1)}L_2p\in S_2,\\
 a_{(4)}p&=L_2(a_{(2)}p)-a_{(2)}L_2p\in K.
\end{align*}
The next mode is scalar and further ones vanish. The current-normal
and derivative portions again use the general identities in the
bootstrap proof. Thus the zero-mode gates really do supply the
all-singular-mode hypothesis of the recursive theorem.

Apply Theorem~\ref{thm:c-recursive} with strong-generator spaces
$K,S_2,S_3,P_4$. The gates, positive-mode bootstrap, and recursive
recipes establish its hypotheses. Therefore the normal span is a
vertex subalgebra. It contains the original generators of $U$, and
every one of its generators is an actual $U$-vector, so it equals $U$.

A normal expression involving only positive-weight generators can have
weight one only if it is one of the weight-one generators themselves.
Thus
\begin{equation}\label{eq:c-properness}
 U_1=K=\mathfrak{sl}_2\oplus\mathfrak{sl}_2,
 \qquad\dim U_1=6<86=\dim B_1.
\end{equation}
In particular $U$ is proper. The same normal-span argument recovers
$U_2=S_2$ and $U_3=S_3$ from the lower identities, giving
\[
 \operatorname{ch}_U(\tau)=q^{-1/3}
       (1+6q+120q^2+956q^3+O(q^4)).
\]
The statement is stronger than agreement of a finite character prefix:
it rules out extra currents generated at any later stage.

\section{What a finite-field calculation proves}
\index{finite-field rank test}

Exact arithmetic over a small number field can still produce large
rational expressions. Reduction modulo a suitable prime is often an
efficient way to certify a nonzero determinant. The elementary
principle is worth separating from more delicate uses.

\begin{proposition}[A nonzero minor survives in characteristic zero]
Let a matrix have entries in an integral domain of characteristic zero,
and let a specialization to a finite field be defined on its entries
and relevant denominators. If an $r\times r$ minor has nonzero image,
that minor was nonzero before specialization. Hence the original rank
is at least $r$.
\end{proposition}
\begin{proof}
The determinant is a polynomial in the entries, so specialization
commutes with it. A zero determinant could not have nonzero image.
\end{proof}

This proves a lower rank, not an upper rank. If the ambient space has
independently known dimension $322$, a rank-$322$ reduced matrix proves
spanning in characteristic zero. Without that dimension bound, the same
computation would not prove exhaustion. Similarly, a residual that is
zero modulo one prime need not vanish in characteristic zero. Membership
identities in this construction are checked as complete number-field
equalities; finite fields select witnesses or certify nonzero minors.
Example~\ref{ex:verification-modular-rank} works through a
two-by-two matrix whose columns become dependent at one admissible
prime and stay independent at another.

For a reducible coefficient algebra, one must also distinguish one
embedding from all algebraic conjugate components. Later irrational-axis
arguments use units or all split fibers to control every component.
This chapter's basic lesson is already sufficient: state the direction
of every rank implication before interpreting a numerical-looking
table as a theorem.

\section{Exercises}

\begin{exercise}\label{ex:c-closure-1}
Prove the negative-mode identity in the final paragraph of the proof of
Theorem~\ref{thm:c-recursive}. Why does it complete vertex closure once
singular products and normal products are controlled?

\par\smallskip\noindent{\footnotesize\normalfont
\hyperref[sol:ex:c-closure-1]{Solution on page~\pageref*{sol:ex:c-closure-1}.}}\par
\end{exercise}
\begin{exercise}\label{ex:c-closure-2}
In a tensor product of irreducible $\mathfrak{sl}_2$ representations
of dimensions $m,n$, prove that the quotient by $E_LP+E_RP$ is
one-dimensional. Give an explicit representative.

\par\smallskip\noindent{\footnotesize\normalfont
\hyperref[sol:ex:c-closure-2]{Solution on page~\pageref*{sol:ex:c-closure-2}.}}\par
\end{exercise}
\begin{exercise}\label{ex:c-closure-3}
Check the constituent count and dimension of the displayed decomposition
of $P$. Recover $\dim(E_LP+E_RP)=294$ without a matrix calculation.

\par\smallskip\noindent{\footnotesize\normalfont
\hyperref[sol:ex:c-closure-3]{Solution on page~\pageref*{sol:ex:c-closure-3}.}}\par
\end{exercise}
\begin{exercise}\label{ex:c-closure-4}
Construct an integer matrix whose rank drops modulo $p$, and explain
why a reduced rank gives only a lower bound on the characteristic-zero
rank. Construct an integer residual that vanishes modulo $p$ but is
not zero.

\par\smallskip\noindent{\footnotesize\normalfont
\hyperref[sol:ex:c-closure-4]{Solution on page~\pageref*{sol:ex:c-closure-4}.}}\par
\end{exercise}
\begin{exercise}\label{ex:c-closure-5}
For a weight-two quasiprimary $a$, derive the $L_1$ commutator used in
the positive-mode bootstrap. Then derive
$2a_{(1)}p=L_1(a_{(0)}p)-a_{(0)}L_1p$.

\par\smallskip\noindent{\footnotesize\normalfont
\hyperref[sol:ex:c-closure-5]{Solution on page~\pageref*{sol:ex:c-closure-5}.}}\par
\end{exercise}
\begin{exercise}\label{ex:c-closure-6}
Explain, in the proof of recursive closure, why induction on total
weight alone is insufficient and why the second lexicographic entry
repairs it. Identify the unique type of term for which that second
entry is needed.

\par\smallskip\noindent{\footnotesize\normalfont
\hyperref[sol:ex:c-closure-6]{Solution on page~\pageref*{sol:ex:c-closure-6}.}}\par
\end{exercise}
